\section{Methodology}
\label{sec:methodology}

\subsection{Prediction target}
\label{sec:methodology-target}

The Primary Confirmatory Endpoint is Target A (Adverse Price Shift) at
depletion threshold $X = 50\%$, forecast horizon $h = 500$\,ms, and
buffer $\Delta = 50$\,ms: the label is positive if the best price level
(bid for a long-side alert, ask for a short-side alert) moves adversely
within $[t+\Delta,\ t+\Delta+h]$. The pre-registered effect-size threshold
is $\Delta\PRAUC \ge +0.005$ over $M_1$ at $\alpha = 0.05$, fixed before
any result was examined. A secondary/exploratory grid additionally varies
the depletion threshold $X \in \{30\%, 50\%, 70\%\}$ and horizon
$h \in \{100\text{ms}, 500\text{ms}, 2\text{s}\}$ (buffer fixed at 50\,ms)
for both Target A and Target B (Pure Depth Depletion: best price
constant, queue size drops $>X\%$ within the same window), for a nominal
$2 \times 3 \times 3 = 18$ specifications
(Section~\ref{sec:results-secondary-grid}).

\subsection{Model ladder}
\label{sec:methodology-ladder}

Four models are compared, each nested in the features available to the
next:

\begin{itemize}[leftmargin=2em]
\item \textbf{$M_0$ (Baseline L1 LOB):} level-1 order-flow imbalance
  $\mathrm{OFI}_{L1}$, micro-price basis $(P_{\text{micro}} - P_{\text{mid}})$,
  bid-ask spread, and absolute queue sizes with rolling historical
  percentile rank $\mathrm{Rank}(Q(t)) \in [0,1]$. Omitting the percentile
  rank term makes $M_0$ artificially weak and inflates the apparent gain
  of every model built on top of it, so it is included by design.

\item \textbf{$M_1$ (Multi-Scale EWMA Bank -- unrestricted tree learning):}
  $M_0$ plus 10-event $\times$ 5-scale exponential event counts,
  \[
    Z_j^{(\beta_p)}(t) = \int_0^t e^{-\beta_p (t-s)}\,dN_j(s), \qquad
    \beta_p \in \{0.2,\ 1.0,\ 5.0,\ 25.0,\ 100.0\}\ \mathrm{s}^{-1}, \quad
    j \in \{1,\dots,10\},
  \]
  fed to an unconstrained LightGBM gradient-boosted-tree model.

\item \textbf{$M_2$ (Parametric Hawkes -- inductive-bias benchmark):}
  $M_0$ plus the MLE-fitted multivariate Hawkes intensity
  \[
    \boldsymbol{\lambda}(t) = \boldsymbol{\mu}(t) + \sum_{p=1}^{5} \mathbf{A}^{(p)} \mathbf{Z}^{(\beta_p)}(t),
  \]
  using the \emph{same} $\beta$ grid as $M_1$, fed as an additional
  feature set to the same LightGBM architecture. As noted in
  Section~\ref{sec:intro-structural-note}, $\boldsymbol{\lambda}(t)$ is
  structurally a linear function of $M_1$'s own feature bank; the MLE fit
  is unsupervised (fits event likelihood, never the forecasting target)
  and can only impose a regularization prior over the same space, not add
  new information.

\item \textbf{$M_3$ (Structural Non-Nested Dynamics):} $M_1$ plus a rolling
  hourly spectral radius $\rho(\boldsymbol{\Gamma}_t)$ (warm-started,
  non-overlapping hourly refits, forward-filled onto the 1\,Hz feature
  grid as a causal step function) and linearized trade-size marks applied
  only to \texttt{trade\_buy}/\texttt{trade\_sell} events. Both blocks are
  genuinely non-linear or externally-covariate-derived summaries that
  $M_1$'s raw event-count EWMA bank does not contain in any decayed-linear
  form, so $M_3$ is the only model in the ladder that could, structurally,
  exceed $M_1$ for a reason other than overfitting noise. The regime
  feature is explicitly an hourly-resolution indicator and is not expected
  a priori to carry signal at the 500\,ms Primary Endpoint horizon.
\end{itemize}

\subsection{Hawkes MLE formulation}
\label{sec:methodology-hawkes}

The full marked, multivariate intensity used for MLE calibration is
\[
  \lambda_i(t) = \mu_i(t) + \sum_{j=1}^{10} \sum_{p=1}^{P}
    \Big( \alpha_{ij}^{(p)} Z_j^{(\beta_p)}(t) + \alpha_{ij}^{\prime(p)} Z_{j,\mathrm{mark}}^{(\beta_p)}(t) \Big),
\]
with trade-size marks $m_k = \mathrm{Volume}_k / \mathrm{Median}(\mathrm{Volume})_{[t-1\mathrm{h},\,t]}$
computed causally (strictly excluding event $k$ itself from its own
normalizing median) and entering \emph{additively} rather than as a
bilinear $(1+\gamma m)\alpha$ product, which preserves global concavity
of the log-likelihood. The baseline intensity follows a parametric
diurnal and asymmetric funding-jump model,
\[
  \mu_i(t) = \Big(a_0 + \sum_{k=1}^{8}\big(a_k\cos(\tfrac{2\pi k t_{\mathrm{sec}}}{86400}) + b_k\sin(\tfrac{2\pi k t_{\mathrm{sec}}}{86400})\big)\Big)
  \cdot \exp\Big(\sum_{f=1}^{3}\big(\delta_{\mathrm{pre}}\mathbb{1}_{[-5\mathrm{m},0)} + \delta_{\mathrm{post}}\mathbb{1}_{[0,+5\mathrm{m}]}\big)\Big),
\]
with funding times $t_{\mathrm{fund}} \in \{00{:}00, 08{:}00, 16{:}00\}$
UTC and asymmetric pre-/post-funding jump coefficients.

The \emph{dimensionless branching-ratio matrix}
\[
  \boldsymbol{\Gamma} = \sum_{p=1}^{5} \mathbf{A}^{(p)} / \beta_p
\]
governs long-run stability: the process is stationary (subcritical) iff
the spectral radius $\rho(\boldsymbol{\Gamma}) < 1$. The specification
calls for enforcing this constraint \emph{during} MLE; the implementation
instead runs an unconstrained per-target L-BFGS-B fit and applies a
post-hoc uniform radial rescaling of all $\alpha$ (and $\alpha'$)
coefficients to $\rho = 0.995$ only if the unconstrained fit exceeds that
ceiling. This deviation from the specification was tested empirically
across 185 fits spanning the full static $N$-grid, both tie rules, and a
purpose-built hourly stress test, and the projection never activated in
that scan (closest approach $\rho = 0.9875$, 0.75\% headroom); it was
later found to activate exactly once, at $N=31$d, in one hourly window
that coincides with the 2024-03-08 \texttt{bookTicker} gap (a known,
already-documented data artifact, not a general large-$N$ property),
affecting 0.116\% of that budget's training rows -- discussed further in
Section~\ref{sec:limitations}.

MLE calibration subsamples a stratified selection of events per fitting
window rather than using the full raw event stream, for L-BFGS-B
wall-clock tractability; the calibration budget and its sensitivity to
sample-size scaling are examined directly in
Section~\ref{sec:results-calibration}.

\subsection{LightGBM setup}
\label{sec:methodology-lgbm}

$M_0$, $M_1$, $M_2$, and $M_3$ are all realized as LightGBM
gradient-boosted-tree classifiers on their respective feature sets. The
headline learning-curve results
(Section~\ref{sec:results-learning-curve}) use a coarse, fixed
3-bucket hyperparameter heuristic applied identically to $M_1$ and $M_2$
at every training-sample budget $N$. A dedicated follow-up
(Section~\ref{sec:results-tuning}) instead tunes \texttt{num\_leaves},
\texttt{min\_data\_in\_leaf} (equivalently \texttt{feature\_fraction} and
\texttt{lambda\_l2} in the tuned search), and related regularization
parameters independently per $N$ and independently for $M_1$ and $M_2$,
via a causal (time-ordered, no shuffling) train/validation split and
random search, specifically to test whether the untuned heuristic was
biasing the $M_1$-vs-$M_2$ comparison in either direction.

\subsection{Evaluation metrics}
\label{sec:methodology-metrics}

Three metrics are reported throughout: area under the precision-recall
curve (PR-AUC, the primary ranking-quality metric, appropriate given
class-imbalanced positive rates), the Brier score (mean squared error
between predicted probability and binary outcome, a calibration-sensitive
metric), and log-loss (binary cross-entropy, also calibration-sensitive
and the metric on which the formal Giacomini--White test is defined).
PR-AUC and Brier/log-loss can and do disagree in this paper's results
(Section~\ref{sec:results-m3}) -- a model can rank events no better while
still producing better-calibrated probabilities, and the two metrics are
not redundant.

\subsection{Giacomini--White test with Newey--West HAC standard errors}
\label{sec:methodology-gw}

The formal significance test used throughout is the Giacomini--White
\citep{giacomini2006tests} conditional predictive ability (CPA) test, run under a genuinely
out-of-sample rolling walk-forward protocol rather than a single static
train/test split. Concretely: (1) compute the per-event loss differential
$d_t = L(y_t, \hat y_t^{M_a}) - L(y_t, \hat y_t^{M_b})$ (log-loss unless
noted otherwise); (2) aggregate into 1-minute non-overlapping mean blocks
$\bar d_\tau$; (3) run the GW CPA test on $\{\bar d_\tau\}$ with
Newey--West/HAC standard errors; (4) re-estimate both models on a
rolling basis -- retrain every 24 hours on the preceding trailing
$N$-day window, forecast the next day -- so that the test's asymptotics,
which require a genuinely out-of-sample rolling forecast sequence, are
valid. The primary rolling GW protocol used in this paper trains on a
trailing 7-day window, re-estimates daily, and forecasts
2024-02-29 through 2024-03-30 (31 out-of-sample days).

\subsection{Romano--Wolf stepdown FWER control}
\label{sec:methodology-rw}

Family-wise error rate across the 18-specification secondary/exploratory
grid (Section~\ref{sec:results-secondary-grid}) is controlled via a
Romano--Wolf stepdown bootstrap \citep{romano2005stepwise,romano2005exact},
resampled by day-blocks to preserve
intraday dependency structure. The Primary Endpoint
(Section~\ref{sec:methodology-target}) is exempt from this correction --
it is tested once, standalone, at $\alpha = 0.05$, per the pre-registered
protocol. Both a cheaper static 7-day-test-week design (7 day-blocks) and
the full 31-day rolling protocol (31 day-blocks, matching the primary GW
protocol above) were run; results are reported for both, with the 31-day
rolling version treated as the more rigorous of the two.

\subsection{Difference-in-differences placebo control}
\label{sec:methodology-placebo}

A placebo design is specified for validating $M_3$'s structural claim:
permute event-type labels $j \in \{1,\dots,10\}$ at fixed timestamps
(preserving timing, destroying type-specific structure), and require
that $M_3$'s \emph{advantage over} $M_1$ (not merely $M_3$'s own
performance) collapse under scrambling, together with a secondary check
lagging point-process feature inputs by $\tau=10$s and verifying that
predictive advantage decays monotonically toward zero. This placebo
control was not run in the present study (see
Section~\ref{sec:limitations}); given the null PR-AUC result for $M_3$
on the Primary Endpoint (Section~\ref{sec:results-m3}), the placebo test
was judged to have little to falsify at the time and was deprioritized
relative to isolating the source of $M_3$'s calibration effect via
ablation, which was run instead.

\subsection{Operational lead-time metric}
\label{sec:methodology-operational}

Independent of the statistical tests above, an operational metric is
reported at a fixed false-alarm budget (top 1\% alert rate, by each
model's own score distribution): the fraction of adverse events captured
(recall at 1\% alert rate) and the median/mean lead time between an
alert and the qualifying adverse event. No execution simulation,
slippage model, or strategy return is computed -- this remains a
forecasting-quality translation, not a trading study.
